% =====================================================================
% Section 4: Method — EvidenceFlow
% =====================================================================
%
% Status: Draft §4 Method v5 — advisor rewrite (2026-05-16)
%   v1: 5 subsections — backed up at archive/04_method_v1_5sections_20260515.tex
%   v2: 3 subsections — collapsed §4.3+§4.4+§4.5 into "Query-Local Evidence Flow"
%   v3: unnumbered overview + 3 mechanism subsections
%   v4: §4.1 rewritten as Offline Indexing (PropRAG-style)
%   v5: full rewrite per advisor feedback —
%        - Method body never references other systems (deferred to §2 Related Work)
%        - §4.3 reframed from reader-boundary to retrieval enhancement
%        - §4.2 explicitly includes local-PPR over query-local subgraph
%        - All "fact-witness" compound terms removed; "OpenIE facts" + "provenance" retained
% Framing reference: research_memory/emnlp_expand_then_compose/46_*.md
% Writing discipline: research_memory/emnlp_expand_then_compose/45_*.md
%
% Method name: \methodname (defined via \newcommand in main file)
%
% Discipline checklist (v5):
%   - No internal codename anywhere in main paper
%   - No reverse-framing sentences ("not X, but Y" against other systems) inside Method
%   - All baseline contrasts deferred to §2 Related Work
%   - Reader budget K is an implementation detail, not a method-level objective
%   - Output is a ranked evidence list R_q over passages, truncatable for any reader
%
% =====================================================================

\section{Method: \methodname{}}
\label{sec:method}

We introduce \methodname{}, a graph-based retrieval framework for multi-hop
question answering. Given a question $q$ and a corpus $\mathcal{D}$,
\methodname{} returns a ranked evidence list $R_q$ over passages in
$\mathcal{D}$. The method has three components: an offline
document-grounded OpenIE graph that supplies entity anchors and relational
structure (\S\ref{sec:method:substrate}), an evidence-driven local graph
retrieval procedure that ranks passages over a question-induced subgraph
(\S\ref{sec:method:retrieval}), and a retrieval-enhancement step that
improves multi-hop coverage on top of the local ranking
(\S\ref{sec:method:enhancement}). The central object underlying these
components is a query-conditioned mass distribution $\pi_q$---which we
refer to as the \emph{evidence flow}---propagated from question anchors
through provenance-grounded edges of the offline graph; \methodname{}
names this propagation together with the two enhancement steps that
operate on top of it.

\subsection{Offline Indexing: Document-Grounded Multi-Layer OpenIE Graph}
\label{sec:method:substrate}

We build a document-grounded multi-layer OpenIE graph
$\mathcal{G}_{\mathcal{D}}=(\mathcal{V}_{\mathcal{D}},\mathcal{E}_{\mathcal{D}})$
over the corpus $\mathcal{D}$ during offline indexing. The graph supplies
the entity anchors and relational structure that \methodname{} draws on at
query time (\S\ref{sec:method:retrieval}). Following standard OpenIE
practice \citep{angeli2015leveraging,gutierrez2025hipporag}, we extract
salient assertions from each document and organize them into the graph
below; the extractor itself is not a contribution of this work.

\paragraph{OpenIE fact extraction.}
For each document $d \in \mathcal{D}$, a language model extracts a set of
OpenIE facts $f = (h, r, t)$, where $h$ and $t$ are head and tail entities
and $r$ is the relation phrase that connects them. We retain each fact
together with its provenance $\mathrm{src}(f) \in \mathcal{D}$, so that
every fact can be traced back to the document from which it was extracted.
The resulting fact set is denoted $\mathcal{F}$.

\paragraph{Graph construction.}
The graph $\mathcal{G}_{\mathcal{D}}$ contains two node types and three
edge types. Nodes are: (i) \emph{passage nodes}, one per document
$d \in \mathcal{D}$, carrying the natural-language content that a
downstream reader will consume; and (ii) \emph{phrase nodes}, one per
distinct entity surface form extracted across $\mathcal{F}$. Edges are:
(i) \emph{passage--phrase} edges that connect each document to the
entities it mentions; (ii) \emph{relational} edges that connect a pair of
phrase nodes $(h, t)$ whenever some fact $f = (h, r, t) \in \mathcal{F}$
asserts a relation between them, with $r$ retained as an edge attribute
and $\mathrm{src}(f)$ retained as the edge provenance; and (iii)
\emph{synonymy} edges that connect phrase nodes whose dense embeddings
are highly similar, providing alias coverage across surface variants of
the same entity.

\paragraph{Indexing for retrieval.}
We pre-compute dense embeddings for all passage and phrase nodes, and
materialize $\mathcal{G}_{\mathcal{D}}$ as an adjacency index keyed by
node identifier. The graph is built once and reused across all questions:
at query time, no new OpenIE extraction or fact generation is performed
over $\mathcal{D}$. The artifacts produced by offline indexing---phrase
nodes, passage--phrase links, relational edges with provenance, and
synonymy edges---are exactly the pieces \methodname{} draws on when it
constructs a query-local subgraph and ranks passages by graph propagation
(\S\ref{sec:method:retrieval}).

\subsection{Evidence-Driven Local Graph Retrieval}
\label{sec:method:retrieval}

Given a question $q$, \methodname{} induces a query-local subgraph
$G_q = (V_q, A_q)$ from $\mathcal{G}_{\mathcal{D}}$ and runs personalized
PageRank locally over $G_q$ to rank passages. The subgraph isolates the
neighborhood of $\mathcal{G}_{\mathcal{D}}$ that is relevant to $q$, and
local propagation lets passage relevance be shaped by entity and
relational connectivity rather than by independent surface similarity to
$q$.

\paragraph{Query-local subgraph construction.}
We extract a question anchor set $A(q) \subseteq \mathcal{E}$ from $q$ by
identifying entity mentions in the question, using a controlled language
model that is applied to the question only, never to retrieved documents.
We also compute an entry signal $\mathrm{Entry}(q) \subseteq \mathcal{D}$
from a standard dense retriever to seed activation. The subgraph $G_q$ is
initialized with (i) phrase nodes corresponding to anchors in $A(q)$, (ii)
passage nodes in $\mathrm{Entry}(q)$, and (iii) every passage node that is
the provenance of at least one relational edge whose endpoints include an
anchor in $A(q)$. Starting from this seed set, we expand $G_q$ along
relational, passage--phrase, and synonymy edges of
$\mathcal{G}_{\mathcal{D}}$ for a bounded number of hops, keeping the
activated region compact relative to $\mathcal{D}$ while still allowing
passages reachable only through multi-hop entity chains to enter $V_q$
even when their surface similarity to $q$ is low.

\paragraph{Local personalized PageRank.}
Within $G_q$, OpenIE facts act as the propagation units of the evidence
flow: two facts $f_i, f_j \in \mathcal{F}$ are connected when they share
an activated entity, and a fact $f$ is linked to its provenance passage
$\mathrm{src}(f)$. The seed distribution $s_q$ places mass on facts whose
arguments include an anchor in $A(q)$, weighted by the entry scores of
their provenance passages and any entry-side mass on $\mathrm{Entry}(q)$.
We solve, by power iteration restricted to $G_q$,
\begin{equation}
\label{eq:local-ppr}
\pi_q = \alpha\, s_q + (1-\alpha)\, P_q^{\!\top}\, \pi_q,
\end{equation}
where $P_q$ is the local transition matrix induced by these fact-fact
and fact-passage links. We then read out a passage score by aggregating
the stationary mass on the facts that ground the passage:
\begin{equation}
\label{eq:passage-score}
s_{\mathrm{ppr}}(d \mid q)
=
|\mathcal{F}_d|^{-1/2}
\sum_{f \in \mathcal{F}_d} \pi_q(f).
\end{equation}
Here, $\mathcal{F}_d = \{\, f \in \mathcal{F} : \mathrm{src}(f) = d \,\}$
is the set of OpenIE facts grounded in passage $d$.
The square-root normalization avoids penalizing passages whose facts have
been deduplicated more aggressively. Because propagation is restricted to
$G_q$, the resulting scores depend on whether a passage is connected to
$q$ through entity-relation evidence in $\mathcal{G}_{\mathcal{D}}$, and
not only on whether it is independently similar to $q$. The output of
local retrieval is the resulting ranked list of passages in $V_q$, which
the next subsection refines.

\subsection{Retrieval Enhancement with Relational Evidence}
\label{sec:method:enhancement}

Local propagation yields an initial graph-aware ranking, but two failure
modes remain. First, passages that lie on a multi-hop chain but are far
from question anchors in $G_q$ can receive low propagated mass, even when
they carry the relational evidence required to complete the chain.
Second, the top of the local ranking can be dominated by passages that
cluster around a single entity neighborhood, leaving other entities or
relations demanded by $q$ under-covered. We address both with a
retrieval-enhancement step that operates on top of the local-PPR scores.

\paragraph{Relation-grounded expansion.}
For each high-scoring node in $\pi_q$, \methodname{} follows the relational
edges of $\mathcal{G}_{\mathcal{D}}$ to retrieve additional passages whose
provenance realizes a relation that has not yet been covered by the
current top of the ranking. Concretely, for a phrase or passage node $u$
with high $\pi_q(u)$, we enumerate the relational edges incident to $u$
whose relation phrase $r$ is compatible with the relations that $q$ asks
about, and admit the provenance passages of those edges as expansion
candidates. Each expansion candidate inherits a score that combines
$\pi_q(u)$ with an edge-level score derived from the dense similarity
between the relation phrase $r$ and the question. This step brings in
bridge passages and tail-entity passages that are required by the
question's relational structure but are not lexically similar to $q$.

\paragraph{Coverage-aware reranking.}
After expansion, we rerank the combined candidate set under a noisy-OR
coverage model over a small set of question-side requirements
$B(q) = \{b_1, b_2, \ldots\}$, where each $b \in B(q)$ is an entity- or
relation-bearing slot that $q$ expects the evidence to jointly resolve.
The requirements are extracted from $q$ by the same controlled language
model used for the anchor set; they are applied to the question only,
never to retrieved documents, and do not generate evidence chains. For a
passage $d$ and a requirement $b$, we associate a contribution
$\phi_{d,b} \in [0,1]$ derived from dense similarity between $d$ and $b$
under $q$. The coverage of a candidate set $R$ for requirement $b$ is
\begin{equation}
\label{eq:noisy-or}
\mathrm{cov}_q(R \mid b)
\;=\;
1 - \prod_{d \in R}\bigl(1 - \phi_{d,b}\bigr),
\end{equation}
which saturates at $1$ once any passage in $R$ contributes substantially
to $b$, and rewards diverse contributions across passages otherwise. We
form $R_q$ greedily by adding, at each step, the candidate $d$ that
maximizes the marginal gain
$\sum_{b \in B(q)}\bigl[\mathrm{cov}_q(R \cup \{d\} \mid b) - \mathrm{cov}_q(R \mid b)\bigr]$
weighted against $s_{\mathrm{ppr}}(d \mid q)$. Because the noisy-OR
coverage saturates per requirement, passages whose neighborhoods are
already covered by higher-ranked candidates receive diminishing marginal
gain, while passages addressing previously-uncovered requirements are
promoted. The reranked output is the final ranked evidence list $R_q$
over passages in $\mathcal{D}$; downstream systems may consume any prefix
of $R_q$.
